\section{研究の意義}

本研究の意義は，以下の三点にまとめられる．

第一に，軟弱地面に固有の鉛直方向の不確実性を明示的に扱う点である．従来の歩行制御では主に水平方向の安定性に注目してきたが，本研究では床沈み込みに起因する重心高さ変動を歩行不安定化の主要因として捉え，鉛直方向の運動補償を導入することで新たな安定化手法を示している．

第二に，LIPMの枠組みを維持したまま重心高さ変動に対応可能である点が挙げられる．提案手法は，複雑な非線形モデルや最適化問題を新たに導入することなく，既存のLIPMベース歩行生成器に対して比較的少ない変更で実装可能である．このため，実装容易性および計算効率の観点から，実機応用への展開が期待できる．

第三に，シミュレーションを通じて床沈み込み環境下での有効性を定量的に示している点である．バネマスダンパ系としてモデル化した軟弱床を用いた飛び石地形シミュレーションにより，従来手法と提案手法の挙動を比較し，重心発散抑制および歩行継続性の向上を確認した．これにより，本研究の手法が軟弱地面に対するロバストな歩行制御手法として有効であることを示している．

以上より，本研究は，二足歩行ロボットの実環境適応性を高めるための基礎的かつ実用的な知見を提供するものであり，将来的な屋外歩行ロボットの実用化に貢献する意義を有すると考えられる．